\chapter{The threat}

Cryptography and device controls stop remote and
algorithmic attack. They do not stop a person with a
wrench, a kidnapping crew, or a detention at a border.
Physical security is the work of protecting people from
that, and of making the keys they hold less useful to
whoever is hurting them.

A \textbf{wrench attack} is any case where the attacker
skips the protocol and goes after the human who can
unlock, sign, or coerce someone else to sign. The
assumption is physical reach, or a credible threat of it.
Phishing and malware are different problems.

The name comes from the xkcd comic in which 4096-bit RSA
loses to a cheap wrench.\footnote{\url{https://xkcd.com/538/}}
The joke landed because it is how a lot of theft actually
works.

\section{Why this is a book now}

More value sits in self-custody, non-custodial wallets,
and on-chain treasuries, often on a handful of people
with no institutional physical security.

Public chains, social media, and leaks make it easier to
tie a real name to a balance.

Wallet documentation still spends its pages on malware
and seed handling. The physical and social path is
treated as someone else's problem. It is not.

This book assumes the attacker can do basic OSINT, watch
a building, and pressure a person. It does not assume a
nation-state. Some of them will still use the same holes.

\section{Who shows up}

\textbf{Opportunistic criminals.} Local crews who learn
that a target ``has crypto'' from lifestyle, talk, or
social media. They want liquid funds fast: hot wallets,
exchange apps, whatever is on the phone. Time and visible
risk work against them.

\textbf{Organized crime.} People who combine OSINT,
surveillance, and sometimes an insider. They will spend
time on a founder or a treasury signer.

\textbf{Insiders and the close circle.} Current or former
staff, contractors, partners, family, or an intimate
partner. They already know where devices live, when
approvals happen, and who else can sign. That knowledge
is the attack.

\textbf{State actors.} Security services, police units,
or proxies in places with weak courts or hostile
enforcement. Detention, a travel stop, or a raid, used to
get keys or force an on-chain action.

\section{What they want}

\textbf{Direct theft.} Unlock the phone, say the seed,
sign the transaction.

\textbf{Extortion.} Keep the threat going. Personal funds
and the org treasury are both in play.

\textbf{Operational coercion.} Pause a contract, whitelist
an address, pass a malicious proposal, leak a signing
policy. The wallet is the tool, not the prize.

\textbf{Intelligence.} Map who signs, how recovery works,
where the hardware lives, so the next attempt is quieter.

\section{Where they hit}

\textbf{Home.} Invasions, building common areas, the
parking garage, threats that name family or roommates.

\textbf{Work.} Treasury, ops, or engineering staff on the
commute, at the office door, or in a lobby that lets
visitors wander. Worse when signing devices leave the
building.

\textbf{Public.} Cafés, airports, hotels, nightlife.
Crypto conferences pack targets into one venue and tell
everyone what they do.

\textbf{Family, friends, staff.} Spouses, children,
drivers, assistants. Indirect access, or leverage.

If one person can move the treasury from the couch, the
attacker does not need a clever exploit. They need that
person, at home, after a public tweet.
